\documentclass[conference]{IEEEtran}
\usepackage{cite}
\usepackage{amsmath,amssymb,amsfonts}
\usepackage{graphicx}
\usepackage{array}
\usepackage{textcomp}
\usepackage{xcolor}
\usepackage{url}
\newcommand{\coast}{CoastWeave}
\newcommand{\state}[1]{\texttt{#1}}
\begin{document}

\title{What Could Help Conservation of the Oceans?\\
An Equity-Constrained, Climate-Robust Portfolio Framework for Coastal Conservation}

\author{}
\maketitle

\begin{abstract}
Ocean conservation is often presented as a menu of individually desirable actions, yet nutrient loading, plastic leakage, habitat loss, fishing pressure, climate exposure, capacity, and social legitimacy are coupled across land and sea. This paper scientifically reframes a broad conservation question as a portfolio-decision problem: how can a coastal seascape select intervention combinations that reduce pressures and protect ecological function while respecting implementation capacity, livelihood safeguards, and climate uncertainty? The paper proposes CoastWeave, a design-only, provenance-aware framework that represents nutrient-source reduction, plastic-leakage prevention, habitat and blue-carbon protection or restoration, area-based fisheries management, observing/data access, and governance-capacity strengthening as distinct but connected intervention classes. A non-compensatory gate prevents a high aggregate ecological score from overriding missing monitoring, capacity, legitimacy, or equity evidence. The research design compares an ecological-only classifier with an equity- and capacity-constrained model, evaluates portfolios over declared climate and socioeconomic scenarios, and reports abstention when evidence is too sparse or non-transferable. The contribution is not a claim that a proposed portfolio has improved an actual coast. It is an auditable method for identifying when a cross-sector conservation portfolio is mechanistically plausible, conditionally robust, blocked by equity or capacity, or unsupported by available evidence.
\end{abstract}

\begin{IEEEkeywords}
ocean conservation, coastal management, eutrophication, plastic pollution, climate resilience, environmental equity, marine protected areas, decision support
\end{IEEEkeywords}

\section{Introduction}

Ocean conservation is a global challenge, but global importance alone does not define an actionable scientific problem. Coastal systems receive watershed-derived nutrients and waste, host fisheries and dense infrastructure, contain habitat with climate and biodiversity value, and are governed by overlapping public, private, community, and international institutions. A proposed conservation action may have a plausible ecological mechanism and still fail because it cannot be monitored, is not legitimate to affected people, lacks implementation capacity, shifts burdens to a dependent group, or performs poorly under a changed climate. Conversely, a socially legitimate initiative may be unable to reduce a material pressure if it is disconnected from the source pathways that drive it.

The broad question ``what could help conservation of the oceans?'' is therefore scientifically reframed here as: \emph{How can an equity-constrained, climate-robust, land-sea intervention portfolio reduce nutrient and plastic pressure while protecting coastal ecosystem function and livelihood safeguards in a specified coastal seascape?} The objective is neither to find one universal intervention nor to create a site-specific policy prescription. It is to construct a disciplined decision architecture that exposes what must be true before a portfolio can be described as ready for further local assessment.

The motivation for a portfolio approach comes from the interaction of multiple pressures. Coastal eutrophication is linked to upstream and coastal nutrient dynamics, and management choices can have different effects across the freshwater-marine continuum \cite{paerl2009,malone2020}. Plastic leakage implicates material production, waste systems, behavior, policy, and community participation rather than only removal of debris already in the sea \cite{kumar2021}. Habitat protection and restoration may generate biodiversity and climate co-benefits, but blue-carbon claims require attention to valuation, governance, and socio-economic conditions \cite{hilmi2021}. Area-based protection can be important, but governance, capacity, fisheries integration, and local perceptions influence whether it delivers intended results \cite{mccay2011,bennett2013,gill2017}. These are not substitutes for one another.

This paper introduces \coast{}, an evidence-constrained portfolio framework. It combines a structured intervention ontology, evidence cards, climate and socioeconomic scenarios, non-compensatory feasibility gates, equity/livelihood safeguards, and an abstention state for missing or non-transferable evidence. The proposed work is computational and design-only. It does not establish protected areas, modify fishing access, change land use, install waste infrastructure, restore habitat, collect community data, or make a site-specific recommendation. Its aim is to make the assumptions and evidence limits behind a future conservation claim inspectable.

\section{Evidence Boundary and Research Gap}

\subsection{Land-sea pressure coupling}

The freshwater-marine continuum is a useful reminder that an ocean result may depend on a terrestrial decision. Paerl argues that nutrient management for eutrophication should consider both nitrogen and phosphorus across connected freshwater and coastal systems \cite{paerl2009}. Malone and Newton synthesize how anthropogenic nutrient inputs interact with other pressures, including development, fishing, and climate-related changes, and identify sustained, integrated monitoring and management as necessary for coastal response \cite{malone2020}. These sources support modeling nutrient-source reduction as a linked land-sea lever; they do not specify a universal target, a local response coefficient, or a ready-to-deploy watershed intervention.

The multistressor character is also evident in the Baltic Sea analysis by Reusch \emph{et al.} \cite{reusch2018}. The paper describes a historically disturbed coastal system in which integrated watershed-sea governance and international coordination have mattered, while climate change can offset management effects. The correct lesson is not that every coast should copy a Baltic strategy. It is that a conservation model using only historical mean conditions can hide the fragility of an intervention ranking.

Plastic leakage requires a different but connected pathway. Kumar \emph{et al.} review the ecosystem-service and development implications of plastic pollution and discuss circularity, policy, community involvement, and life-cycle thinking \cite{kumar2021}. The framework therefore treats plastic leakage prevention as an independent portfolio class. It should not be converted into the unsupported claim that one clean-up, ban, or recycling program will solve coastal plastic pollution.

\subsection{Habitat, protection, and social-ecological conditions}

Blue-carbon ecosystems offer an example of a potentially valuable co-benefit rather than an automatic solution. Hilmi \emph{et al.} discuss conservation, protection, and restoration alongside valuation, governance, policy action, and coordination challenges \cite{hilmi2021}. The article was additionally checked through the requested in-app browser at its publisher page: title, journal, DOI, article type, and abstract framing were visible. The evidence supports including habitat/blue-carbon protection or restoration as a conditional portfolio component. It does not support using blue carbon to substitute for emissions mitigation, local participation, pollution control, or ecosystem-specific evidence.

Marine protected areas and other spatial measures must likewise be interpreted through governance. McCay and Jones describe MPAs as place-based, ecosystem-oriented tools that require integration with fisheries management, spatial planning, institutions, and stakeholder engagement \cite{mccay2011}. Bennett and Dearden show that perceived livelihood effects, governance, and management shape local support in a specific Thai MPA context \cite{bennett2013}. Gill \emph{et al.} identify capacity shortfalls as a global hindrance to MPA performance \cite{gill2017}. The collective implication is important: a polygon or designation is an input to a conservation system, not evidence of an outcome.

\begin{table*}[t]
\caption{Evidence classes and permitted use in the proposed portfolio framework.}
\label{tab:evidence}
\centering
\renewcommand{\arraystretch}{1.15}
\begin{tabular}{|p{0.17\textwidth}|p{0.20\textwidth}|p{0.24\textwidth}|p{0.28\textwidth}|}
\hline
\textbf{Evidence class} & \textbf{Representative sources} & \textbf{Permitted use} & \textbf{Use not permitted} \\
\hline
Land-sea nutrient coupling & \cite{paerl2009,malone2020,reusch2018} & Link watershed pressure to coastal portfolio logic; represent climate-sensitive uncertainty & Infer a universal local nutrient target or effect size \\
\hline
Plastic and circularity & \cite{kumar2021} & Represent prevention, waste-system, policy, and community levers as a distinct pressure pathway & Assume a single marine clean-up or ban solves leakage \\
\hline
Habitat and blue carbon & \cite{hilmi2021} & Include conditional habitat/climate co-benefits and governance requirements & Treat habitat action as a substitute for emissions cuts or local consent \\
\hline
MPA governance and capacity & \cite{mccay2011,bennett2013,gill2017} & Gate spatial protection on capacity, legitimacy, and fisheries context & Equate designation with ecological or livelihood benefit \\
\hline
Equity and observing & \cite{bennettsocial2017,blasiak2020,levin2019} & Require participation, distributional safeguards, data access, and observing feasibility & Let an aggregate score erase rights, access, or uncertainty \\
\hline
\end{tabular}
\end{table*}

\subsection{Equity, observing, and scientific legitimacy}

Conservation decisions are affected by who can contribute knowledge, bear restrictions, access benefits, and verify conditions. Bennett \emph{et al.} outline how conservation social science contributes to legitimate and effective environmental action \cite{bennettsocial2017}. Blasiak \emph{et al.} identify unequal capacities in an ocean-knowledge setting \cite{blasiak2020}, and Levin \emph{et al.} emphasize integration, standards, accessibility, and capacity needs for ocean observing \cite{levin2019}. These studies do not offer a ready fairness metric for every coast. They show why a model that reports only aggregate ecological benefit is incomplete.

Five research gaps follow. \emph{G1} is the cross-pressure portfolio gap: evidence is arranged by intervention type rather than by linked choices. \emph{G2} is the land-to-sea decision gap: watershed source pressures and coastal outcomes cross planning boundaries. \emph{G3} is the capacity and legitimacy gap: implementation conditions are routinely secondary to ecological optimization. \emph{G4} is the equity and distribution gap: average gains may hide concentrated livelihood or participation burdens. \emph{G5} is the climate-robustness gap: static rankings can fail when runoff, exposure, ecological state, or socioeconomic capacity shifts. \coast{} addresses these gaps as a research design, not a completed management program.

\section{CoastWeave Framework}

\subsection{Research question and hypothesis}

The primary research question is whether a climate-robust, equity- and capacity-constrained portfolio model produces a materially different conservation classification from an ecological-only model. The central hypothesis is that it will. A portfolio selected only for ecological pressure reduction may look favorable while lacking monitoring feasibility, institution capacity, procedural legitimacy, or livelihood safeguards. A gated model should expose those failures and, where evidence is adequate, select combinations that maintain acceptable conditional performance over stated scenarios.

The framework uses six intervention classes: nutrient-source reduction; plastic-leakage prevention; habitat/blue-carbon protection or restoration; area-based fisheries management; observing and data access; and governance-capacity strengthening. These are categories of analysis, not instructions to alter a coast. Each class has a different evidence mechanism, temporal horizon, and equity implication. A portfolio may include several classes, but a high nominal contribution in one cannot cancel an adverse or unsupported condition in another critical domain.

\begin{figure}[t]
\centering
\footnotesize
\begin{tabular}{c}
\textbf{CoastWeave portfolio decision flow}\\[2pt]
\fbox{\parbox{0.82\columnwidth}{\centering Coastal-seascape evidence package: land-sea pressures, ecosystem state, communities, climate, source limits}}\\[3pt]
$\downarrow$\\[-1pt]
\fbox{\parbox{0.82\columnwidth}{\centering Intervention classes: nutrient reduction; plastic prevention; habitat/blue carbon; area-based fisheries; observing; governance/capacity}}\\[3pt]
$\downarrow$\\[-1pt]
\fbox{\parbox{0.82\columnwidth}{\centering Climate and socioeconomic scenario comparison}}\\[3pt]
$\downarrow$\\[-1pt]
\fbox{\parbox{0.82\columnwidth}{\centering Non-compensatory gates: ecological safeguard, observability, capacity, legitimacy, equity/livelihood}}\\[3pt]
\begin{tabular}{c@{\qquad}c}
\fbox{\parbox{0.36\columnwidth}{\centering Fail or unknown\\Blocked or evidence-abstention state}} &
\fbox{\parbox{0.36\columnwidth}{\centering Pass with sensitivity analysis\\Conditional portfolio-readiness result}}
\end{tabular}
\end{tabular}
\caption{Conceptual CoastWeave flow. The editable diagram source is retained with the four-agent artifacts. This is a proposed decision process, not an observed portfolio evaluation.}
\label{fig:coastweave}
\end{figure}

\subsection{Evidence cards and gates}

The analytical unit is an \emph{Evidence Card}. Each card records a planning scale, ecosystem type, pressure or intervention mechanism, outcome concept, observation horizon, evidence setting, source identifier, uncertainty class, transferability qualifier, and permitted-claim ceiling. A global synthesis can support a general mechanism card, for example, but cannot automatically populate a local effect-size card. A regionally specific study can inform a contextual analogue, but its transfer assumptions remain explicit.

For portfolio $p$, scenario $s$, and critical domain $d$, let $E_{psd}$ be an evidence status and $T_{psd}$ a transferability qualifier. Critical eligibility is defined as

\begin{equation}
R_{ps}=\prod_{d\in\mathcal{D}_{\mathrm{critical}}}\mathbb{I}\left(E_{psd}=\mathrm{sufficient}\right)\mathbb{I}\left(T_{psd}\geq\tau_d\right),
\label{eq:gate}
\end{equation}

where $\mathcal{D}_{\mathrm{critical}}$ includes pressure-mechanism plausibility, ecological safeguard, observability, implementation capacity, procedural legitimacy, and equity/livelihood protection; $\tau_d$ is declared before comparison. The product form is deliberate. If a required domain is adverse, unknown, or non-transferable, $R_{ps}$ prevents the result from being called portfolio-ready regardless of a favorable aggregate score.

For eligible cases only, the study may use a transparent conditional comparison objective,

\begin{equation}
\begin{aligned}
U_{ps}={}&w_EB_{ps}+w_PP_{ps}+w_HH_{ps}\\
&+w_LL_{ps}-w_CC_{ps}, \qquad \sum w=1,
\end{aligned}
\label{eq:objective}
\end{equation}

where $B_{ps}$ denotes ecological-benefit evidence, $P_{ps}$ pressure-reduction evidence, $H_{ps}$ habitat/connectivity evidence, $L_{ps}$ livelihood-safeguard evidence, and $C_{ps}$ conditional implementation burden. Equation \eqref{eq:objective} is not an outcome forecast or a compensatory escape from Eq. \eqref{eq:gate}. It is a fully reported comparison aid among portfolios that have already satisfied critical conditions.

\section{Design-Only Methods}

\subsection{Study structure}

The proposed study does not evaluate a real named coast. It creates a reusable analysis protocol for a future declared planning cell. The planning cell contains linked modules for watershed pressures, coastal habitats, fisheries and connectivity, waste leakage, climate exposure, monitoring/data access, governance capacity, and livelihood/equity safeguards. The model's role is to organize available evidence and make its limitations visible, not to generate a policy decision from abstract literature.

Model A is an ecological-only comparator. It represents nutrient, plastic, habitat, and fisheries/ecosystem proxies without enforcing capacity, legitimacy, monitoring, or distributional requirements. Model B is \coast{}. It applies all modules, the critical gates in Eq. \eqref{eq:gate}, and declared scenarios. The primary output is not a single winning portfolio. It is a traceable comparison that identifies which portfolios change class, which domain caused the change, and whether the change is evidence-supported or assumption-driven.

The initial intervention ontology is intentionally high level. Nutrient-source reduction includes upstream or coastal source-control categories without prescribing instruments. Plastic-leakage prevention includes prevention and circularity categories without specifying municipal operations. Habitat/blue-carbon action includes protection or restoration categories without proposing physical restoration. Area-based fisheries management includes spatial-management categories without defining zones or restrictions. Observing/data access and governance-capacity strengthening are included as interventions because a portfolio may be biologically plausible but non-evaluable or non-implementable without them.

\subsection{Scenario and robustness design}

Each candidate portfolio is evaluated over a bounded scenario matrix. Climate/runoff scenarios modify the plausibility and uncertainty of land-sea nutrient connections, exposure, and habitat stress. Coastal-exposure scenarios represent differences in disturbance and ecological sensitivity without asserting a local forecast. Socioeconomic/implementation scenarios vary institutional capacity, resource constraints, and ability to sustain monitoring. Data-availability scenarios test whether a model's recommendation collapses when the observing inputs it assumes are not present.

Robustness is evaluated as the minimum conditional score across eligible scenarios. A regret quantity records the difference between a portfolio and the highest scoring eligible portfolio in each scenario. These outputs are descriptive, scenario-dependent, and accompanied by sensitivity analysis. No portfolio is called robust merely because it performs well in a historical-average scenario. If a conclusion changes under reasonable threshold, weight, or transferability changes, the report marks it as conditional. If the analysis is dominated by absent evidence or arbitrary assumptions, the output is \state{EVIDENCE\_OR\_MODEL\_ABSTAIN}.

\begin{table}[t]
\caption{Planned CoastWeave decision states.}
\label{tab:states}
\centering
\scriptsize
\renewcommand{\arraystretch}{1.1}
\begin{tabular}{|p{0.25\columnwidth}|p{0.28\columnwidth}|p{0.31\columnwidth}|}
\hline
\textbf{Evidence situation} & \textbf{Permitted state} & \textbf{Required language} \\
\hline
Pressure-intervention mechanism has bounded support, but no local result & Mechanistically plausible & ``Supports a rationale for further evidence assembly, not an observed coastal benefit.'' \\
\hline
Critical gates pass across stated scenarios & Scenario conditional & ``Supports a conditional portfolio-readiness inference under declared assumptions.'' \\
\hline
Ecological score favorable, but capacity, legitimacy, or livelihood gate fails & Equity or capacity blocked & ``Aggregate ecological attractiveness cannot justify a readiness claim.'' \\
\hline
Critical evidence is missing, non-transferable, or assumption-dominated & Evidence/model abstention & ``The available record cannot support a directional local portfolio recommendation.'' \\
\hline
\end{tabular}
\end{table}

\subsection{Stress tests and validation logic}

Four challenge cases test whether the framework behaves as intended. First, an area-based protection or habitat portfolio may appear ecologically favorable while enforcement, monitoring, and governance capacity are inadequate. The gated model should block an implementation-ready interpretation. Second, a nutrient and plastic portfolio may improve an aggregate proxy while shifting cost, access, or participation burdens to a resource-dependent group; the equity/livelihood gate should prevent a favorable classification. Third, a historical optimum may become fragile under a high-runoff or high-exposure scenario; the climate-robust ranking should change or be reported as unstable. Fourth, a model may use data layers that a data-poor coastline cannot sustain; the observability gate should trigger abstention rather than a false precision result.

Quality control begins with provenance. All claims must refer to an Evidence Card and retain its scale, source type, and transfer limitation. A second review checks whether any aggregate score is being used to override a failed gate. A third review checks whether governance, consent, participation, local/Indigenous knowledge, and equity have been reduced to afterthoughts. The framework does not collect or represent local/Indigenous knowledge in this proposal; it records that their absence is a material limitation for any future local assessment.

\section{Conditional Results Logic and Governance Requirements}

The proposed research will not produce observed changes in nutrient loads, plastic leakage, habitat condition, fish populations, livelihoods, or climate risk. Its results are methodological and conditional. A useful result is a reproducible finding that Model B changes the classification obtained from Model A for a stated reason. For example, a biological-only model may favor a habitat or MPA component, while the gated model may classify the portfolio as blocked because implementation capacity is not evidenced. Another useful result is a portfolio whose apparent advantage disappears when its data-access or climate assumptions are relaxed. These would be findings about the quality of a decision claim, not measurements of a coast.

Three outcomes are especially informative. A \emph{pressure-coupling gap} occurs when treating nutrient, plastics, and habitat as one generic benefit obscures a source pathway or trade-off. A \emph{implementation gap} occurs when ecological attractiveness is not matched by monitoring, capacity, legitimacy, or enforcement conditions. A \emph{distribution gap} occurs when an aggregate benefit masks a concentrated burden or inadequate participation safeguard. The presence of these gaps does not imply that conservation is impossible; it identifies the next evidence or co-design task required before a stronger conclusion is legitimate.

Human review is required before any real-world use. Coastal ecologists and hydrologists must assess pressure and habitat pathways. Fisheries, waste-systems, and climate experts must check the intervention and scenario ontology. Governance scholars, social scientists, and legally recognized local/Indigenous knowledge holders must judge whether procedural legitimacy and equity variables are appropriate to the place. This review cannot be simulated by the model because the framework deliberately does not collect stakeholder views, legal information, or local rights data.

The model also avoids a misleading form of optimism: it does not describe blue-carbon, MPA, plastic, or nutrient actions as universal climate solutions. Blue-carbon potential comes with governance and valuation conditions \cite{hilmi2021}; MPAs require capacity and integration \cite{mccay2011,gill2017}; nutrient responses are connected across land and sea \cite{paerl2009,malone2020}; and conservation legitimacy depends on social dimensions \cite{bennettsocial2017}. A portfolio is credible only when these claims are kept distinct while their dependencies are made explicit.

\section{Discussion and Conclusion}

The most important contribution of \coast{} is not a new ecological prediction. It is a stricter language for asking what could help ocean conservation. The answer cannot be a single intervention, a global ranking, or a promise that technology will repair every pressure. It must be a place-sensitive combination of pressure reduction, ecosystem stewardship, observation, capacity, legitimacy, and distributional safeguards. That combination will differ across coasts, but the need to make assumptions, evidence limits, and trade-offs visible is general.

The framework has three advantages. First, it connects land and sea without pretending that all evidence is at one scale. Second, it makes capacity and equity non-compensatory conditions rather than optional evaluation criteria. Third, it tests climate robustness through scenario-dependent minimum performance and regret rather than presenting a historical-optimal solution as durable. These features encourage ambition while preventing a conservation model from overclaiming certainty.

There are material limitations. The evidence synthesis is targeted, not systematic. The framework does not contain local datasets, effect sizes, costs, law, or stakeholder preferences. The intervention ontology is broad and must be co-designed before implementation. Weight choices in Eq. \eqref{eq:objective} are normative as well as technical; they must be declared and sensitivity-tested. A high abstention rate may be scientifically appropriate in data-poor or socially underrepresented settings, but it is not itself a solution to inequity. Finally, a decision framework cannot substitute for emissions reductions, pollution governance, ecological monitoring, or local political legitimacy.

Ocean conservation can nevertheless be approached as a solvable design challenge in the limited sense intended here. Existing evidence supports integrated nutrient management, plastic prevention, habitat protection, spatial governance, observing, and capacity building as relevant classes of action. The next scientific task is to determine which combinations remain defensible when climate uncertainty, implementation limits, and fairness constraints are visible. \coast{} offers a proposal-level method to make that determination auditable and falsifiable before real interventions are considered.

\appendices
\section{Scenario Protocol and Falsification Matrix}

The framework is intentionally designed to be falsifiable at the level of decision logic. Before a future planning-cell analysis begins, investigators must publish a scenario register, the critical-gate thresholds, the intervention ontology, and the definition of every proxy. The register distinguishes inputs observed for the planning cell from literature-derived mechanism assumptions and from values included only for sensitivity analysis. No value may silently change category after a portfolio has been ranked. This requirement is particularly important in coastal systems because runoff, coastal exposure, livelihoods, and institutional capacity can vary at finer scales than a published synthesis.

The baseline scenario represents a declared reference condition, but it is never treated as a prediction. The high-runoff scenario increases the uncertainty or stress assigned to land--sea nutrient connections and related habitat exposure. The high-coastal-exposure scenario tests a condition in which disturbance, warming, or other climate-sensitive pressures weaken an assumed habitat response. The constrained-capacity scenario reduces the feasibility of sustained enforcement, maintenance, and monitoring. The limited-observability scenario removes or degrades critical data streams. Finally, the distributional-safeguard scenario tests whether an apparent ecological gain depends on unexamined concentration of costs, access loss, or inadequate procedural participation. These are stress-test classes, not numeric claims about a particular coastline.

The comparison is rejected when it fails one of the predeclared falsification tests. A portfolio logic is falsified as climate-robust if its class changes from ready to blocked or abstain under an admissible climate scenario without a documented reason. It is falsified as equity-ready if an ecological aggregate stays favorable while an equity/livelihood gate fails. It is falsified as transferable if a conclusion depends on a source whose scale, ecosystem setting, or observation horizon is outside the card's declared limit. It is falsified as observable if a favorable result survives only while unavailable data are assumed. In each case, the appropriate output is a diagnostic and an evidence request, not an adjusted score that conceals the failure.

\begin{table*}[t]
\caption{Predeclared scenario protocol and falsification tests for a future CoastWeave analysis. These are analytic checks, not simulated or observed coastal outcomes.}
\label{tab:scenario}
\centering
\footnotesize
\renewcommand{\arraystretch}{1.14}
\begin{tabular}{|p{0.17\textwidth}|p{0.24\textwidth}|p{0.23\textwidth}|p{0.23\textwidth}|}
\hline
\textbf{Stress-test class} & \textbf{What is varied} & \textbf{Decision failure to detect} & \textbf{Permitted response} \\
\hline
Climate/runoff & Land--sea pressure linkage, habitat stress, and uncertainty bounds & Baseline ranking reverses or critical ecological evidence loses transferability & Report scenario instability; request pathway-specific evidence \\
\hline
Coastal exposure & Disturbance sensitivity and climate-related ecosystem exposure & Claimed habitat or connectivity benefit depends on an untested benign exposure assumption & Block robust-readiness claim; retain conditional rationale only \\
\hline
Capacity and observability & Enforcement, maintenance, monitoring continuity, and data availability & A portfolio remains favorable only because unverified capability or data are assumed & Apply capacity/observability gate or evidence/model abstention \\
\hline
Equity and livelihoods & Participation safeguard, burden distribution, and access-risk variables & Ecological aggregate masks a concentrated burden or weak procedural legitimacy & Apply non-compensatory equity gate; identify co-design and review need \\
\hline
Transferability & Scale, ecosystem analogue, observation horizon, and source setting & Claim exceeds the permitted ceiling on its Evidence Card & Downgrade claim to mechanism or analogue; prohibit local-effect statement \\
\hline
\end{tabular}
\end{table*}

This protocol supports a stronger research claim than generic ``integration.'' It asks whether a proposed portfolio remains defensible after the assumptions that normally sit outside ecological scoring are surfaced and challenged. A result that fails is useful: it identifies precisely whether the next task is ecological measurement, watershed analysis, monitoring investment, governance assessment, or rights-aware co-design. The research therefore advances conservation preparation by identifying decision fragility before a real intervention is proposed.

\section{Evidence Coverage and Human-Review Route}

The evidence plan uses five linked coverage questions. First, what pressure pathway is being addressed, and where does its mechanism apply? Second, what ecosystem function or safeguard is expected to change, and on what observation horizon? Third, what capacity and observing conditions are required to make the claim testable? Fourth, whose participation, rights, access, and livelihood conditions may be affected by the framing of the portfolio? Fifth, how does climate uncertainty alter each of the preceding answers? The five questions prevent the model from treating an evidence-rich ecological component as a substitute for missing social, governance, or measurement evidence.

The required review route begins with a trace review. A marine ecologist, hydrologist, fisheries scientist, waste-systems specialist, or another appropriate domain expert checks the pressure--mechanism statements and flags unsupported causal language. A climate and observing reviewer then checks scenario bounds, data dependencies, and whether the proposed proxy is observable at the stated planning scale. A governance and social-science review examines whether capacity, accountability, procedural fairness, and distributional consequences have been reduced to a token variable. Finally, legally recognized local and Indigenous knowledge holders must be able to determine whether a future local use has represented context, rights, and knowledge appropriately. This proposal does not substitute a literature-derived variable for those judgments.

The reviews are sequential only for traceability; they are not a claim that one discipline can veto the scientific content of another. A failed or unresolved review produces a named uncertainty, an abstention, or a bounded conditional statement. The report must preserve disagreement rather than averaging it away. For example, a biologically plausible habitat component with strong literature support may remain in a portfolio ontology while its readiness state is blocked until capacity and rights-related questions are properly assessed. Conversely, a community-prioritized option with limited transferable ecological evidence may motivate evidence assembly without being represented as a proven ecological intervention.

\begin{table}[t]
\caption{Minimum evidence coverage and review record for each portfolio component.}
\label{tab:review}
\centering
\scriptsize
\renewcommand{\arraystretch}{1.12}
\begin{tabular}{|p{0.18\columnwidth}|p{0.24\columnwidth}|p{0.20\columnwidth}|p{0.16\columnwidth}|}
\hline
\textbf{Coverage domain} & \textbf{Required trace record} & \textbf{Reviewer competence} & \textbf{Failure state} \\
\hline
Pressure and ecosystem pathway & Source, mechanism, scale, horizon, and claim ceiling & Coastal ecology, hydrology, fisheries, or pollution expertise & Mechanism unsupported or non-transferable \\
\hline
Climate and observation & Scenario assumption, proxy, data provenance, and sensitivity result & Climate and observing expertise & Scenario-fragile or not observable \\
\hline
Capacity and governance & Implementation dependency, accountability, and monitoring continuity & Governance and policy expertise & Capacity or legitimacy blocked \\
\hline
Equity and livelihoods & Participation condition, distributional risk, and rights-related limitation & Social science plus appropriate local/Indigenous review & Equity or livelihood blocked \\
\hline
Synthesis integrity & Gate status, score use, abstentions, and dissent record & Cross-disciplinary methods review & Evidence or model abstention \\
\hline
\end{tabular}
\end{table}

The route is designed to make an ambitious conservation programme more credible, not to stall it. A portfolio can proceed from a general rationale to a bounded, evidence-assembly agenda when the necessary gate has not yet passed. That distinction avoids the false choice between universal prescriptions and inaction. It channels uncertainty toward concrete, reviewable research tasks while keeping claims proportional to what the record supports.

\section{Evidence Card and Claim-Validation Contract}

Every \coast{} assessment must publish an Evidence Card registry. A card contains the source and DOI; the planning and ecosystem scale; the pressure/intervention mechanism; the outcome concept; the evidence setting; the observation horizon; the uncertainty class; transferability limits; and a permitted-claim ceiling. The ceiling is the strongest statement the card supports without hidden extrapolation. A global review can support a mechanism statement but not a site-specific outcome. A regional case can support an analogue but not a universal policy conclusion. A local result, if later available, must still state its measurement window and social context.

Each critical domain is assigned one of four statuses: sufficient, adverse, insufficient, or not applicable. ``Sufficient'' means adequate for a bounded decision state, not that environmental risk is absent. ``Adverse'' means the available record identifies a barrier incompatible with the state. ``Insufficient'' produces an abstention when the domain is critical. ``Not applicable'' requires a documented rationale because apparent irrelevance can conceal a missing safety, justice, or governance question.

Claim validation has three passes. The scientific pass verifies source-to-claim fit, including scale and horizon. The methods pass verifies that no score overrides a failed gate and that scenario/weight assumptions are reported. The governance and equity pass verifies that capacity, rights, participation, burden distribution, and knowledge access have not been omitted after ecological benefits are calculated. Any disagreement is retained in the trace record. The final report must list all abstentions and what evidence would be needed to replace them.

This proposal records explicitly what it did not do: it did not sample ecosystems, collect social data, modify waste systems, impose spatial restrictions, restore habitat, or observe an intervention outcome. The contract separates literature evidence from new observations so that a future study can improve a single domain without retrospectively inflating the certainty of other domains.

An auditable assessment also needs a versioned decision record. For every portfolio--scenario comparison, the record must preserve the evidence-card identifiers, gate status by domain, scenario register version, transferability qualifiers, weights, missingness treatment, reviewer comments, and the exact language issued by the classifier. A later local measurement may strengthen one card, but it must not rewrite the interpretation produced from an earlier record. If a threshold or ontology category changes, the analysis is rerun as a new version and the earlier conditional result remains visible. This makes it possible to distinguish an improved evidence base from a retrospective change in decision logic.

The record must contain negative as well as favorable findings. It must list components excluded by a failed gate; portfolios that moved from favorable to blocked under climate or capacity stress; assumptions that dominated a ranking; and domains where the correct output was abstention. A report that preserves only the best scoring combination would not test the framework's central premise. The scientifically useful output is a map of decision fragility: which evidence gap changes a conclusion, which condition cannot be compensated by an aggregate benefit, and which expert review is needed before a local question can be responsibly posed.

For reproducibility, the future implementation should publish machine-readable card and scenario schemas together with a human-readable claim ledger. The ledger links each sentence-level conclusion to the relevant gate, cards, uncertainty qualifier, and reviewer status. It does not replace domain expertise, but it enables a reviewer to locate the source of a strong statement and to test whether that statement survives a declared counterfactual. In this way, CoastWeave treats transparent revision as a research feature rather than as a weakness: a portfolio becomes more decision-ready by resolving named uncertainties, not by concealing them.

\begin{thebibliography}{00}
\bibitem{paerl2009} H. W. Paerl, ``Controlling Eutrophication along the Freshwater-Marine Continuum: Dual Nutrient (N and P) Reductions are Essential,'' \emph{Estuaries and Coasts}, vol. 32, no. 4, pp. 593--601, 2009, doi: 10.1007/s12237-009-9158-8.
\bibitem{malone2020} T. C. Malone and A. Newton, ``The Globalization of Cultural Eutrophication in the Coastal Ocean: Causes and Consequences,'' \emph{Frontiers in Marine Science}, vol. 7, 670, 2020, doi: 10.3389/fmars.2020.00670.
\bibitem{reusch2018} T. B. H. Reusch \emph{et al.}, ``The Baltic Sea as a time machine for the future coastal ocean,'' \emph{Science Advances}, vol. 4, no. 5, eaar8195, 2018, doi: 10.1126/sciadv.aar8195.
\bibitem{kumar2021} R. Kumar \emph{et al.}, ``Impacts of Plastic Pollution on Ecosystem Services, Sustainable Development Goals, and Need to Focus on Circular Economy and Policy Interventions,'' \emph{Sustainability}, vol. 13, no. 17, 9963, 2021, doi: 10.3390/su13179963.
\bibitem{hilmi2021} N. Hilmi \emph{et al.}, ``The Role of Blue Carbon in Climate Change Mitigation and Carbon Stock Conservation,'' \emph{Frontiers in Climate}, vol. 3, 710546, 2021, doi: 10.3389/fclim.2021.710546.
\bibitem{mccay2011} B. J. McCay and P. Jones, ``Marine Protected Areas and the Governance of Marine Ecosystems and Fisheries,'' \emph{Conservation Biology}, vol. 25, no. 6, pp. 1130--1133, 2011, doi: 10.1111/j.1523-1739.2011.01771.x.
\bibitem{bennett2013} N. J. Bennett and P. Dearden, ``Why local people do not support conservation: Community perceptions of marine protected area livelihood impacts, governance and management in Thailand,'' \emph{Marine Policy}, vol. 44, pp. 107--116, 2013, doi: 10.1016/j.marpol.2013.08.017.
\bibitem{gill2017} D. A. Gill \emph{et al.}, ``Capacity shortfalls hinder the performance of marine protected areas globally,'' \emph{Nature}, vol. 543, pp. 665--669, 2017, doi: 10.1038/nature21708.
\bibitem{bennettsocial2017} N. J. Bennett \emph{et al.}, ``Conservation social science: Understanding and integrating human dimensions to improve conservation,'' \emph{Biological Conservation}, vol. 205, pp. 93--108, 2017, doi: 10.1016/j.biocon.2016.10.006.
\bibitem{blasiak2020} R. Blasiak \emph{et al.}, ``The ocean genome and future prospects for conservation and equity,'' \emph{Nature Sustainability}, vol. 3, no. 8, pp. 588--596, 2020, doi: 10.1038/s41893-020-0522-9.
\bibitem{levin2019} L. A. Levin \emph{et al.}, ``Global Observing Needs in the Deep Ocean,'' \emph{Frontiers in Marine Science}, vol. 6, 241, 2019, doi: 10.3389/fmars.2019.00241.
\bibitem{duarte2020} C. M. Duarte \emph{et al.}, ``Rebuilding marine life,'' \emph{Nature}, vol. 580, pp. 39--51, 2020, doi: 10.1038/s41586-020-2146-7.
\end{thebibliography}
\end{document}
